\documentclass[11pt]{article}
\usepackage{enumitem}
\usepackage[margin=1in]{geometry}
\usepackage{amsmath}
\usepackage{booktabs}
\usepackage{graphicx}
\usepackage{hyperref}
\usepackage{titling}
\setlength{\parskip}{3pt}
% paragraph spacing — replaces plain \usepackage{parskip}
\usepackage[skip=3pt plus 1pt minus 0.5pt]{parskip}

% heading spacing
\usepackage{titlesec}
\titlespacing{\section}      {0pt}{1.6ex plus .2ex}{0.8ex}
\titlespacing{\subsection}   {0pt}{1.4ex plus .2ex}{0.6ex}
\titlespacing*{\paragraph}    {0pt}{1.0ex plus .2ex}{1em}
\usepackage[
    backend=biber,
    style=numeric,
    sorting=none,
    doi=true,
    url=true,
    eprint=true
]{biblatex}

\addbibresource{references.bib}

% Move title block upwards
\setlength{\droptitle}{-6em}

\title{Deep Learning for Multivariate Time Series Forecasting}

\author{
\small
Group ID: 75 \\
\small
Nikolas Martin Diehl (2797049) \\
\small
Ceylin Ekinci (3429785) \\
\small
Kshitij Parashar (3440735) \\
\small
Fabian Joel Zwinge (3758126)
}

\small\date{\today}

\begin{document}

\maketitle
\vspace{-0.6em}

% ============================================================
% Abstract
% ============================================================

\begin{abstract}
\noindent Multivariate time-series forecasting underlies many operational decisions. We study a benchmark in
which an hourly operational load index is forecast across 96 series over a 336-hour horizon, with
eleven domain covariates \emph{known at inference} while the target is withheld. We compare
gradient-boosted trees, a gated multilayer perceptron (MLP), and a Temporal Fusion Transformer
(TFT). The TFT attains the best validation WAPE (14.48), substantially improving on the naive
last-value (48.10) and seasonal-mean (34.41) baselines, but only narrowly beating the gated
MLP (14.75) and the LightGBM-MLP ensemble (14.68). This suggests a feature-set ceiling in which
different architectures extract similar information. On a cow tri-axial accelerometer dataset,
the same TFT architecture beats the naive baseline by roughly 32\% (WAPE 43.0 vs.\ 63.8);
an ablation shows a modest benefit from the known behaviour covariate.
\end{abstract}


% ============================================================
% 1. Introduction
% ============================================================

\section{Introduction}
Multivariate time-series forecasting is relevant in operational settings where
decisions depend on accurate long-horizon forecasts. In this project, we forecast
an hourly operational load index across 96 time series over a horizon of 336 hours.
Eleven future covariates are known at inference time, while the future target remains
hidden. The task therefore differs from history-only forecasting and can be viewed
as known-covariate sequence-to-sequence forecasting.
\newline 

\noindent We investigate this problem across several model classes. Gradient-boosted
decision trees~\cite{lightgbm} serve as traditional machine-learning baselines, while
a gated multilayer perceptron (MLP) represents a generic deep-learning approach with
learned feature selection. The central model is the Temporal Fusion Transformer
(TFT)~\cite{tft}, designed for multi-horizon forecasting with static, historical,
and known future inputs. Comparing these approaches allows us to examine whether
performance is driven mainly by model architecture or by the available feature set.
We further investigate whether the findings transfer to a smaller cow accelerometer
dataset~\cite{cowdata}.
\newline

\noindent Our research questions are:

\begin{itemize}
    \item \textbf{RQ1.}
    Can a deep-learning model outperform strong traditional machine-learning
    baselines (e.g.\ gradient-boosted trees) on this
    known-future-covariate task?

    \item \textbf{RQ2. (core)}
    Does an architecture purpose-built for known future covariates with
    variable selection (TFT) beat a generic gated MLP - or does a feature-set
    ceiling hold across architecture classes, with diverse models converging
    to the same function on fixed features?

    \item \textbf{RQ3.}
    Do the findings transfer to a second forecasting dataset (the cow
    accelerometer data, cast into the same setup), and does augmentation help
    more in that low-data regime?
\end{itemize}


% ============================================================
% 2. Related Work
% ============================================================

\section{Related Work}
Time-series forecasting includes classical methods such as ARIMA, ETS, and
seasonal averages, as well as machine-learning and deep-learning approaches.
Gradient-boosted trees, such as LightGBM, can model nonlinear relationships but
typically depend on manually engineered temporal features~\cite{lightgbm}.
Recurrent models such as LSTMs learn sequential dependencies directly from
historical observations~\cite{lstm}. More recent architectures include
PatchTST~\cite{patchtst}, iTransformer~\cite{itransformer}, and pretrained
foundation models such as Chronos~\cite{chronos}. Qiu et
al.~\cite{channelsurvey} provide a broader overview of deep multivariate
forecasting from a channel-strategy perspective.
\newline

\noindent Model complexity does not necessarily lead to better forecasts. DLinear, for
example, demonstrates that simple linear models can compete with transformer
architectures on long-horizon tasks~\cite{dlinear}. The TFT addresses multi-horizon forecasting with heterogeneous inputs by
separately processing static, historical, and known future variables. It also
uses adaptive variable-selection networks and supports probabilistic forecasts
through quantile outputs~\cite{tft}. Another relevant research direction is
time-series data augmentation, which includes techniques such as window
sampling, perturbation, masking, and frequency-domain transformations and is
particularly relevant in settings with limited training data~\cite{augsurvey}.


% ============================================================
% 3. Method
% ============================================================

\section{Method}


% ------------------------------------------------------------
% 3.1 Input representation
% ------------------------------------------------------------

\subsection{Input Representation and Feature Engineering}

Before defining the final model inputs, we performed an exploratory data analysis (EDA) of the target, covariates, temporal patterns, correlations, missing values, and differences between the individual time series. Based on these observations and subsequent experiments, the final input representation combines historical target information, known future covariates, static series information, and engineered temporal features.
\newline

\noindent
The eleven provided domain covariates are available to both the encoder and the full forecast horizon. In addition, first-order differences of \texttt{queue\_pressure\_forecast}, \texttt{network\_pressure\_forecast}, and \texttt{shock\_risk} are included to capture short-term changes. Temporal information is represented through cyclical encodings of hour of day, day of week, and day of year, together with categorical hour and weekday features. The historical target is used only by the encoder, while \texttt{series\_id} identifies the individual time series.
\newline

\noindent
Missing covariate values are imputed separately within each series using forward filling followed by backward filling. The target is normalized per series using a group-based normalizer, while the training loss is computed on the original target scale.
\newline

% ------------------------------------------------------------
% 3.2 Model Selection
% ------------------------------------------------------------
\subsection{Model Selection and Baselines}

Following the initial data analysis and feature engineering, we explored several forecasting approaches to identify suitable model families. We first established statistical and machine-learning baselines, including LightGBM~\cite{lightgbm}, and then evaluated deep-learning approaches such as a gated MLP and LSTM encoder-decoder models~\cite{lstm}, PatchTST~\cite{patchtst}, and Chronos~\cite{chronos}.
\newline

\noindent Some simpler approaches provided strong reference models and were retained as baselines for our research questions. In contrast, architectures such as PatchTST and Chronos were less promising in our initial setup and were not pursued further. These experiments narrowed the model space and provided the reference points against which our final forecasting architecture is evaluated.

% ------------------------------------------------------------
% 3.2 Architecture
% ------------------------------------------------------------
\subsection{Model Architecture}
Our final model is based on the TFT, which combines recurrent sequence processing, variable selection, gating, and attention to model both short- and long-term dependencies in multivariate time series. Figure~\ref{fig:architecture} illustrates the overall architecture. It consists of five main building blocks:

\begin{figure}[h]
    \centering
    \includegraphics[width=0.85\textwidth]{tft.png}
    \caption{Overview of the TFT forecasting architecture. Figure from original paper~\cite{tft}.}
    \label{fig:architecture}
\end{figure}

\begin{enumerate}
    \item \textbf{Gating Mechanisms.}
    Gating mechanisms control the information flow through the network and allow unnecessary transformations to be suppressed or skipped.

    \item \textbf{Variable Selection Networks.}
    Variable selection networks determine which historical observations and known future covariates are most relevant at each time step.

    \item \textbf{Static Covariate Encoding.}
    Static covariates are encoded into context vectors that can condition other parts of the network.
    These representations provide additional context for variable selection and temporal processing.

    \item \textbf{Temporal Processing.}
    An LSTM-based sequence-to-sequence structure captures local and short-term temporal dependencies from historical observations and known future inputs.
    The resulting representations are subsequently processed by interpretable multi-head attention to capture longer-term relationships across time.

    \item \textbf{Forecast Generation.}
    The model uses the P50 quantile, which represents the median and provides a single point forecast for each future time step.
    As the forecasting task requires one predicted value per timestamp, the model produces 336 P50 forecasts, which are subsequently evaluated using the different performance metrics considered in our experiments.
\end{enumerate}

\noindent The final TFT uses a hidden size of 96, a hidden continuous size of 32, four attention heads, one LSTM layer, and a dropout rate of 0.15.
% ------------------------------------------------------------
% 3.3 Training
% ------------------------------------------------------------

\subsection{Training Setup}
Training samples use an encoder history of up to 336 hours and a fixed prediction horizon of 336 hours. To increase training-window diversity, the encoder length is randomly varied between 168 and 336 hours, with windows sampled at a stride of 24 hours. Model selection uses three-fold rolling-origin validation while preserving temporal order. Each fold evaluates a 336-hour period, while the preceding 336 hours are used exclusively for early stopping, ensuring that evaluation data do not influence model selection.
\newline

\noindent
The model is optimized using MAE on the original target scale and AdamW with a learning rate of $10^{-3}$ and weight decay of $10^{-4}$. We use a batch size of 32 with gradient accumulation over four batches, resulting in an effective batch size of 128. Training runs for up to 60 epochs with an early-stopping patience of 10, gradient clipping at 0.5, and dropout of 0.15. After model selection, the final models are trained on the full training period for 10 epochs, corresponding to the median best epoch from cross-validation. Five models with seeds 42, 7, 13, 99, and 1234 are trained and their predictions averaged to form a five-model ensemble. All models are trained using mixed precision on two NVIDIA Tesla T4 GPUs.

% ============================================================
% 4. Experiments
% ============================================================

\section{Experiments}


% ------------------------------------------------------------
% 4.1 Experimental Setup
% ------------------------------------------------------------

\subsection{Experimental Setup}
The TFT model's performance is evaluated using WAPE, MAE, MSE, RMSE, MAPE, and sMAPE. It is compared against last-value, seasonal-mean, LightGBM, and gated-MLP baselines.

% ------------------------------------------------------------
% 4.2 Validation Results
% ------------------------------------------------------------

\subsection{Validation Results}

Table~\ref{tab:benchmark-results} compares the forecasting approaches on the
benchmark validation set.

\begin{table}[h]
\centering
\begin{tabular}{lrrrr}
\toprule
Method & WAPE & MAE & RMSE & Notes \\
\midrule
Naive last-value baseline & 48.10 & 5.29 & 6.97 & Minimum baseline \\
Seasonal mean baseline & 34.41 & 3.79 & 5.21 & Stronger reference \\
LightGBM & 14.68 & 1.62 & 3.06 & Machine-learning baseline \\
Gated MLP & 14.75 & 1.62 & 3.06 & Main reference model \\
TFT & 14.48 & 1.59 & 3.08 & Final architecture \\
\bottomrule
\end{tabular}
\caption{Validation results on the benchmark dataset. Lower is better for all listed metrics.}
\label{tab:benchmark-results}
\end{table}

\noindent
The TFT achieves a WAPE of 14.48 and an MAE of 1.59, substantially outperforming the naive last-value baseline (WAPE 48.10) and the seasonal-mean baseline (34.41). This corresponds to a WAPE reduction of approximately 70\% and 58\%, respectively. The gated MLP reaches a WAPE of 14.75, while the best LightGBM-MLP ensemble achieves 14.68.
\newline

\noindent
Although the TFT therefore performs best in terms of WAPE and MAE, the differences between the learned models are small. Their RMSE values are also very similar, with the TFT at 3.08 compared with 3.06 for both alternatives. This suggests that much of the achievable performance is driven by the informative known future covariates, while the TFT's variable selection and temporal modeling provide an additional but relatively modest improvement.
\newline



% ------------------------------------------------------------
% 4.3 Ablation Study
% ------------------------------------------------------------

\subsection{Analysis: Architecture versus Feature-Set Ceiling}

The near-identical scores of the three learned models in Table~\ref{tab:benchmark-results} - TFT (14.48), the
LightGBM-MLP ensemble (14.68), and the gated MLP (14.75), all within roughly $0.3$ WAPE and with
RMSE differing by $0.02$ - act as a de-facto architecture ablation. Despite very different inductive
biases, the models converge to almost the same predictive function once they have access to the same
known future covariates. This is consistent with a \emph{feature-set (information) ceiling} rather
than an architecture ceiling: the eleven domain forecasts largely determine what is predictable, and
the purpose-built known-covariate structure of the TFT yields only a small, though consistent, edge
(RQ2). Residual diagnostics on the shipped model support this reading; the remaining error is
concentrated in a spike tail (the largest-error rows dominate the WAPE numerator) and in per-series
heterogeneity (per-series WAPE varies several-fold), neither of which is removed by changing
architecture. To isolate the contribution of the known-future covariate directly - rather than
inferring it from cross-model convergence - we run a controlled with/without-covariate ablation on
the additional dataset in Section~\ref{sec:cow} (Table~\ref{tab:cow}).
% ------------------------------------------------------------
% 4.4 Augmentation Ablation
% ------------------------------------------------------------
\subsection{Augmentation Ablation}
\label{sec:augmentation}

Since the models in Table~\ref{tab:benchmark-results} converge to nearly the same WAPE regardless
of architecture, we tested whether training-time data augmentation could push the TFT past this
ceiling. This experiment was applied only to the benchmark dataset. Two techniques were validated at the feature level
before training: covariate dropout (randomly withholding a domain covariate per training sample,
$p{=}0.10$) and covariate jitter (Gaussian noise sized at $\sigma{=}0.20\times$ each covariate's
own standard deviation). 

\begin{table}[h]
\centering
\begin{tabular}{lr}
\toprule
Method & WAPE  \\
\midrule
TFT (no augmentation) & 15.27 \\
TFT + covariate jitter ($\sigma{=}0.20\times$std) & 15.84 \\
TFT + covariate drop & 	15.54 \\
\bottomrule
\end{tabular}
\caption{Augmentation ablation. Evaluation coverage differed across runs; the baseline-matched dropout result provides the strongest evidence.}
\label{tab:augmentation}
\end{table}

\noindent Covariate dropout, evaluated under the full 3-fold, 5-seed protocol, was consistently
worse than its matched baseline across all fold averages and all 5 seeds of the final fold.
Covariate jitter showed the same overall trend. A separate single-fold, single-seed baseline is
reported only for reference, since it is more sensitive to run-to-run variance. Overall, neither
augmentation method improved on the unaugmented TFT. This is consistent with a feature-set ceiling:
if performance is mainly limited by the informativeness of the eleven domain covariates, perturbing
them during training is unlikely to provide substantial additional gains.
% ------------------------------------------------------------
% 4.5 Additional Dataset
% ------------------------------------------------------------

\subsection{Additional Dataset}
\label{sec:cow}

To evaluate whether the findings generalize beyond the benchmark, we apply the same architecture to
the Japanese Black Beef Cow Behavior dataset~\cite{cowdata}: six cows recorded with tri-axial neck
accelerometers at 25\,Hz, together with hand-annotated behaviour labels.

\paragraph{Casting the data into the benchmark setup.}
We treat each contiguous cow recording segment as one series (34 in total). The target is
resampled \emph{movement intensity}, defined as the within-second standard deviation of the
signal-vector magnitude
$\text{SMV}=\sqrt{a_x^2+a_y^2+a_z^2}$.
The known-future covariate is the behaviour label, interpreted as a planned schedule, together
with cyclical calendar features. We forecast 30 seconds from 60--120 seconds of history, using
a 6-second stride and holding out the final 30 seconds of each series for validation.

\paragraph{Adaptations.}
Only dataset-specific preprocessing and temporal scales were changed: accelerometer data were
resampled from milliseconds to seconds, the encoder and forecast horizon were rescaled accordingly,
and the behaviour label replaced the benchmark's domain covariates. We used batch size 64 and
three seeds (42, 7, 13).

During research and development the decision was made to refer the impact of augmentation on the additional dataset as stated in RQ3 to future research. The proven impact of augmentation on the original dataset implies potential for optimization in that regime.

\begin{table}[h]
\centering
\begin{tabular}{lrrr}
\toprule
Method & WAPE & MAE & RMSE \\
\midrule
Naive last-value baseline & 63.8 & 0.043 & 0.091 \\
Series-mean baseline & 85.2 & 0.058 & 0.093 \\
TFT, without behaviour covariate (ablation) & 44.5 & 0.030 & 0.059 \\
\textbf{TFT (with behaviour covariate)} & \textbf{43.0} & \textbf{0.029} & \textbf{0.056} \\
\bottomrule
\end{tabular}
\caption{Validation results on the additional (cow) dataset, three-seed mean.}
\label{tab:cow}
\end{table}

\paragraph{Results and discussion.} Table~\ref{tab:cow} reports the transfer results. The TFT reaches
WAPE~43.0, beating the naive last-value baseline (63.8) by about 32\% and the series-mean baseline
(85.2) by half, comfortably clearing the assignment's minimum bar and confirming that the
architecture transfers to a very different, low-data, high-frequency domain. Removing the behaviour
covariate (ablation) worsens WAPE to 44.5 and is also consistently worse on MAE and RMSE, so the
known-future covariate contributes a real gain. That gain is \emph{modest}, however
($\approx$3\% relative WAPE): because the encoder already observes the recent activity signal, target
autocorrelation carries most of the predictable structure, and the covariate adds a smaller increment
on top. This mirrors the benchmark's information-ceiling finding (RQ3): once the dominant signal is
captured, additional structure yields diminishing returns. We did not enable data augmentation in the
reported cow runs; quantifying whether the subsequence-crop augmentation used on the benchmark helps
more in this lower-data regime is left as future work.

% ============================================================
% 6. Conclusion
% ============================================================

\section{Conclusion}

Framing the benchmark around its known-future covariates was the decisive step. A TFT that
architecturally separates the known-future pathway gave our best validation result (WAPE 14.48).
Answering the research questions explicitly:
\begin{itemize}[nosep,leftmargin=1.5em]
  \item \textbf{RQ1 (yes, with a caveat):} the learned models - the TFT, the gated MLP, and the
  gradient-boosted-tree pipeline - all vastly outperform the naive and seasonal baselines, and the
  deep TFT attains the best WAPE/MAE; but its margin over the strong tabular baselines is small.
  \item \textbf{RQ2 (largely a ceiling):} the purpose-built TFT beats the generic gated MLP only
  marginally (14.48 vs.\ 14.75), and all three learned models share almost identical RMSE - evidence
  that the fixed known-covariate feature set, not the architecture, sets the performance ceiling.
  \item \textbf{RQ3 (transfers):} on the cow dataset the TFT again beats the naive baseline
  substantially, and the known-covariate ablation again shows a modest but consistent contribution
  from the known future covariate - the same information-ceiling pattern.
\end{itemize}

\paragraph{Limitations.} Benchmark error is concentrated in a spike tail and in per-series
heterogeneity that a single global model averages over. The transfer dataset is short, high-frequency,
and recorded on a single day, so its absolute scores are a proof of concept rather than a tuned
system; the cow runs used three seeds and no augmentation, and a benchmark-side feature-removal
ablation was not run.

\paragraph{Future work.} Per-series or mixture models to address heterogeneity; a distributional loss
to model the spike tail inside the network; a controlled data-augmentation study in the low-data cow
regime; and a formal benchmark ablation over the known-covariate feature groups.
% ============================================================
% Contributions
% ============================================================

\section*{Contributions}
All four members jointly shaped the project: we scoped the problem together, each proposed a candidate
additional dataset and prototyped model approaches, and then selected as a team both the additional
dataset (Ekinci's proposal) and the single strongest model to develop further.
\newline
\begin{itemize}[nosep,leftmargin=1.5em]
  \item \textbf{Zwinge} and \textbf{Parashar} did the bulk of the work on the core forecasting model:
  reframing the task around the known-future covariates and building the
  LSTM~$\to$~gradient-boosted-tree~$\to$~TFT progression, the cross-validation harness, and the final
  submission archive.
  \item \textbf{Ekinci} initiated the TFT-model approach and proposed the additional dataset chosen by the group (the cow accelerometer
  data), supported the further model development, and pursued the evaluation and analysis of the results.
  \item \textbf{Diehl} led the ideation and wrote the exposé, implemented and adapted the additional-dataset (cow) transfer experiment,
  and - together with Ekinci and Zwinge - polished this report.
\end{itemize}

% ============================================================
% AI Tool-Usage
% ============================================================
\section*{AI Tool-Usage}
AI tools were used throughout each project phase. Claude Code was used for initial exploratory
work with the given data, accelerating the approaches developed by team members. Claude Code and,
additionally, OpenAI's Codex were used to develop the selected TFT approach and adapt it to the
additional dataset, combining vibe-coding with directed bug-fixing. Claude Code and Claude Chat
were used to cross-check references and polish the written results, exposé, and final report for
structure, completeness, and comprehensiveness.


% ============================================================
% References
% ============================================================
\printbibliography

\end{document}